\section{Síntesis y ampliación}

\subsection{Un hilo conductor técnico claro}

Al revisar los casi 20 años de carrera profesional de Chen Tianqi (陈天奇), un hilo conductor lo atraviesa por completo: \textbf{hacer que los sistemas de aprendizaje automático sean más eficientes, más fáciles de usar y más fáciles de desplegar}. Cada proyecto fue una respuesta al cuello de botella del proyecto anterior:

\begin{itemize}
    \item Al participar en el KDD Cup descubrió que necesitaba modelos de árboles eficientes $\rightarrow$ \textbf{XGBoost}
    \item Al trabajar en XGBoost reavivó su entusiasmo por el aprendizaje profundo $\rightarrow$ \textbf{CXXNet/MXNet}
    \item Al trabajar en MXNet descubrió que la mano de obra de ingeniería era el cuello de botella $\rightarrow$ \textbf{TVM} (sustituir los kernels escritos a mano por técnicas de compilación)
    \item TVM necesitaba validarse en un dominio vertical $\rightarrow$ \textbf{MLC LLM} (despliegue de modelos grandes en dispositivos de borde)
    \item La necesidad de comercializar TVM $\rightarrow$ \textbf{OctoML} (código abierto $\rightarrow$ empresa emergente $\rightarrow$ adquirida por NVIDIA)
\end{itemize}

\begin{table}[H]
\centering
\begin{tabular}{p{2.2cm}p{1.8cm}p{9cm}}
\toprule
\textbf{Proyecto} & \textbf{Periodo} & \textbf{Problema que resolvió} \\
\midrule
XGBoost & 2014 & Permitió que los modelos de árboles se ejecutaran más rápido y fueran más fáciles de usar en datos a gran escala (manejo automático de missing value) \\
MXNet & 2015 & Permitió que el marco de trabajo de aprendizaje profundo diera soporte a Python first, programación automática, paralelismo con múltiples GPU y diferenciación automática \\
TVM & 2017 & Permitió que los modelos de ML se adaptaran automáticamente a distintos tipos de hardware (GPU/TPU/NPU), reduciendo los kernels escritos a mano \\
MLC LLM & 2023 & Permitió que los modelos de lenguaje grandes se ejecutaran en dispositivos de borde como teléfonos, navegadores y MacBook \\
\bottomrule
\end{tabular}
\caption{Cronología de las contribuciones técnicas de Chen Tianqi}
\end{table}

La lógica subyacente de este hilo conductor es: \textbf{los sistemas de ML no deberían ser tan complejos; deberían permitir que menos personas los construyeran rápidamente con la técnica correcta}. Esta convicción se mantiene constante desde la era de XGBoost (una sola persona llevándolo al extremo), pasando por la era de TVM (el compilador genera automáticamente los kernels), hasta la era de MLC LLM (despliegue multiplataforma en dispositivos de borde).

\subsection{Lecciones para los investigadores jóvenes}

Con base en el contenido de la entrevista, los consejos centrales que Chen Tianqi da a los jóvenes pueden resumirse en los siguientes puntos, cada uno respaldado por su experiencia personal:

\begin{enumerate}
    \item \textbf{Elegir el problema correcto}: «Trabajar en un problema mediocre y en uno importante toma casi el mismo tiempo.» Esto se lo enseñó el profesor Ao Ping (敖平) de la Universidad Jiao Tong, y lo ha confirmado una y otra vez a lo largo de 20 años. ImageNet fue el problema correcto y las RBM fueron el método equivocado: la elección del problema importa más que la elección del método.

    \item \textbf{No fijar al mismo tiempo el problema y el método}: esta es la «lección más valiosa» que aprendió tras dos años sin resultados trabajando en ImageNet en la Universidad Jiao Tong. Una vez elegido el problema, el método debe ajustarse con flexibilidad: más adelante, el cambio decidido de las RBM a la factorización de matrices en el KDD Cup fue la aplicación de esta lección.

    \item \textbf{Llevarlo al extremo}: la exigencia de Carlos era «write the best paper in the world». No se trataba de perseguir la cantidad de artículos, sino de llevar cada proyecto al nivel de best paper. El éxito de XGBoost es el mejor ejemplo de «llevar un algoritmo al extremo».

    \item \textbf{Aceptar el fracaso}: «Trabajar en un problema muy difícil implica necesariamente incertidumbre y fracaso.» Fracasar temprano es, en realidad, algo bueno: te vuelve más capaz de actuar con libertad y evita la dependencia de la trayectoria.

    \item \textbf{Seguir el instinto}: al elegir, pregúntate «si fracaso, ¿me arrepentiré?». No es un análisis racional de costo-beneficio, sino una manera de asegurarte de que sigues el camino ideal.

    \item \textbf{Partir de first principle}: no temer entrar en campos desconocidos. Desde escribir un transpiler hasta construir un compilador y diseñar un NPU, Chen Tianqi siempre se lanzó de inmediato, como «un ternero recién nacido que no le teme al tigre». La capacidad de diseñar sistemas se afina con la práctica.

    \item \textbf{El código abierto es una oportunidad única para los estudiantes universitarios}: en una empresa empiezas por un componente pequeño dentro de un proyecto grande; en un proyecto de código abierto tienes la oportunidad de diseñar la arquitectura de un sistema desde cero. Este tipo de entrenamiento es difícil de conseguir en otro lugar.
\end{enumerate}

\begin{knowledgebox}{Un consejo especial para los estudiantes de doctorado}
Cuando le preguntaron: «si hoy fueras un estudiante de doctorado de tercer año y solo pudieras elegir corregir o cambiar una cosa, ¿qué harías?», la respuesta de Chen Tianqi fue: «Creo que elegiría seguir haciendo lo que hago ahora.» Es decir, continuar con la compilación de aprendizaje automático: permitir que menos personas, con menos engineering resource, resuelvan el problema del bring up de los sistemas de ML en el hardware más reciente. No cambiaría de rumbo solo porque la época es distinta, porque este problema «está lejos de resolverse».
\end{knowledgebox}

\subsection{Perspectivas sobre el futuro de la IA}

A Chen Tianqi no le interesa definir qué es la AGI («no me interesa demasiado cómo se defina la AGI en sí»); le importa más que sea «pragmatically útil»: ¿puede un equipo de una sola persona generar la producción de un equipo de 100 personas? La tecnología actual «sí nos permite avanzar más rápido hacia esa meta»: por ejemplo, él mismo usa Cursor para refactorizar TVM, y sin herramientas de IA «refactorizar sería mucho más difícil».

Respecto a cómo se configurará el mundo de la IA en el futuro, espera ver que florezcan cien flores: que el código abierto y los individuos también puedan innovar, en lugar de que unos cuantos gigantes dominen todo. Cuando le preguntaron si «la probabilidad de que aparezca ese mundo está aumentando o disminuyendo», su respuesta fue cautelosa: «Solo puedo decir que requiere el esfuerzo de cada persona. Sin duda, distintas personas tendrán creencias distintas.»

La línea técnica en la que insiste siempre ha sido: \textbf{hacer más con menos recursos y de manera más automática}. «Tal vez se resuelva en unos años, tal vez no se resuelva necesariamente, pero al menos es un problema muy interesante.»

Esta frase quizá sea el mejor resumen de toda la carrera de Chen Tianqi: no persigue el éxito garantizado, sino los problemas interesantes. Mientras un problema sea suficientemente importante e interesante, vale la pena invertir en él: incluso si el resultado es el fracaso, también es «un proceso de volver a partir».

\subsection{Conclusión}

Al final de la entrevista, el conductor dejó un espacio abierto para que Chen Tianqi hablara libremente. Dijo:

\textbf{«A lo largo de todo el proceso he tenido mucha suerte de recibir la ayuda de muchísimas personas. Lo que más profundamente me ha dejado es esto: todavía podemos aceptar el fracaso. En el proceso de aceptar el fracaso, hay que seguir animándose a hacer cosas interesantes. Cada proceso es, en realidad, un proceso de volver a partir. We need to reinvent ourselves, then keep doing the next thing.»}

Si tuviera que decirle una frase a su yo del pasado y otra a su yo del futuro, su elección no sería mirar hacia atrás desde el presente para enseñarle algo a su yo pasado, sino al contrario: necesita que aquel yo pasado, «un ternero recién nacido que no le teme al tigre», le recuerde a su yo presente: \textbf{recordar el compromiso consigo mismo, mantener el ideal y seguir adelante.}

El conductor resumió la entrevista diciendo: el Chen Tianqi ideal de dentro de 10 o 20 años debería ser «un Chen Tianqi con más fracasos, capaz de conservar esa mentalidad original para enfrentar lo que queremos enfrentar». Chen Tianqi estuvo de acuerdo con gusto: «Habrá más fracasos, y podré seguir conservando esa mentalidad original para enfrentarlos. Así sería lo mejor.»

\begin{importantbox}{Frases clave de la entrevista}
A continuación se presentan algunas de las frases originales más reveladoras de esta entrevista:
\begin{itemize}
    \item «Trabajar en un problema bastante mediocre y en uno importante toma casi el mismo tiempo.»
    \item «El tiempo es limitado, así que hay que hacer lo que uno realmente quiere hacer.»
    \item «Muchas veces tu valor es, en realidad, un compromiso que tu yo pasado hizo contigo mismo.»
    \item «Si eliges este camino y fracasas, ¿te arrepentirías?»
    \item «El fracaso tampoco es tan terrible. Hay que aceptar el hecho de que se puede fracasar.»
    \item «Creo que lo más difícil para mí ahora sigue siendo mantener la intención original.»
    \item «Hacer tecnología excelente siempre es acertado.»
    \item ``We need to reinvent ourselves, then keep doing the next thing.''
\end{itemize}
\end{importantbox}

\subsection{Lecturas adicionales}

\begin{itemize}
    \item Chen Tianqi, «Diez años de investigación en aprendizaje automático», Zhihu, 2019. Este artículo repasa su trayectoria de investigación de 2009 a 2019, y muchas de las historias que aparecen ahí se desarrollan con más detalle en esta entrevista. La frase del artículo «un proceso aleatorio siempre converge a un estado estable proporcional al esfuerzo invertido» se difundió ampliamente.
    \item Artículo de XGBoost: Chen \& Guestrin, ``XGBoost: A Scalable Tree Boosting System'', KDD 2016. Este artículo tiene hasta la fecha casi 70,000 citas y es uno de los artículos más citados en la historia del KDD.
    \item Artículo de MXNet: Chen et al., ``MXNet: A Flexible and Efficient Machine Learning Library for Heterogeneous Distributed Systems'', NeurIPS Workshop 2015. Los autores aparecen en orden alfabético, lo cual refleja el espíritu de «colaboración pura» entre estudiantes de doctorado.
    \item Artículo de TVM: Chen et al., ``TVM: An Automated End-to-End Optimizing Compiler for Deep Learning'', OSDI 2018. Dio origen a la línea de investigación de la compilación de aprendizaje automático.
    \item Artículo de LIME: Ribeiro, Singh, Guestrin, ``Why Should I Trust You? Explaining the Predictions of Any Classifier'', KDD 2016. Otro trabajo pionero del laboratorio de Carlos del mismo año, que abrió la línea de la interpretabilidad en el aprendizaje automático.
    \item Proyecto MLC LLM: \url{https://github.com/mlc-ai/mlc-llm}. Compila y despliega modelos de lenguaje grandes en distintos dispositivos de borde.
    \item Proyecto Apache TVM: \url{https://tvm.apache.org/}. Infraestructura de código abierto para la compilación de aprendizaje automático.
    \item Curso MLC (el curso de compilación de aprendizaje automático que Chen Tianqi imparte en CMU): \url{https://mlc.ai/}. Explica de manera sistemática los principios y la práctica de la compilación de aprendizaje automático.
    \item XGrammar: \url{https://github.com/mlc-ai/xgrammar}. Biblioteca de código abierto para la generación de datos estructurados, que se ha convertido en uno de los estándares de facto para la salida estructurada de agentes de LLM.
    \item Segundo episodio del podcast WhynotTV (entrevista a Hu Yuanming (胡渊鸣)): también aborda temas como la comercialización de proyectos de código abierto y el emprendimiento académico, y ofrece un contraste interesante con este episodio.
\end{itemize}

\end{document}
